\begin{slide}{}
\begin{center}
	An Architecture for Declarative\\
	Programming\\[1in]
	Dissertation Defense\\
	Michael L. Hilton\\[1in]
\end{center}
\begin{flushleft}
\underline{GOAL}: Develop a computing architecture to implement the 
{\em complete} $\lambda\beta$-calculus efficiently.
\end{flushleft}
\end{slide}

\begin{slide}{}
\begin{center}
	\underline{Outline}\\[.6in]
\end{center}
\begin{flushleft}
\begin{itemize}
\item $\lambda\beta$-Calculus
\item Explain Implementation Difficulties
\item Present RED2 Solutions
\item Assessment of the RED2 Architecture
\item Enhancements \& Expansions 
\end{itemize}
\end{flushleft}
\end{slide}



%3
\begin{slide}{}
\begin{center}
	The Untyped $\lambda\beta$-Calculus\\[.75in]
\end{center}

\underline{SYNTAX}

$x \in $ Identifiers\\
$e \in \lambda$-Expressions

$e$ ::=  $x \mid \lambda x . e \mid e_1 e_2$\\[.5in]

\underline{SEMANTICS}

$\alpha$-Conversion:
$\lambda x . e\  \longleftrightarrow\ \lambda y .[y/x]e$

$\beta$-Reduction:
$(\lambda x . e_1) e_2$ \red $[e_2/x]e_1$

\begin{flushleft}
Where ``$[e_2/x]e_1$'' means ``substitute $e_2$ for every occurrence of
$x$ in $e_1$''
\end{flushleft}
\end{slide}

\begin{slide}{}
\begin{center}
SUBSTITUTION
\end{center}

1. $[e_1/x]x = e_1$\\
2. $[e_1/x]y = y$, where $x \not= y$\\
3. $[e_1/x](e_2 \ e_3) = ([e_1/x]e_2 \ [e_1/x]e_3)$\\
4. $[e_1/x](\lambda x.e_2) = \lambda x.e_2$\\
5. $[e_1/x](\lambda y.e_2) = \lambda y . [e_1/x]e_2$,\\ \hspace*{2em}
		if $x \not= y$ and \\ \hspace*{2em} either $y \not\in fv(e_1)$ 
		or $x \not\in fv(e_2)$\\
6. $[e_1/x](\lambda y.e_2) = \lambda z . [e_1/x]([z/y]e_2)$,\\ \hspace*{2em}
		if $x \not= y$ and $y \in fv(e_1)$ and $x \in fv(e_2)$,\\ \hspace*{2em}
		and $z \not\in fv(e_1) \cup fv(e_2)$

\begin{center}
FREE VARIABLE
\end{center}

1. $fv(x) = \{x\}$\\
2. $fv(e_1 \ e_2) = fv(e_1) \cup fv(e_2)$\\
3. $fv(\lambda x.e) = fv(e)-\{x\}$
\end{slide}


\begin{slide}{}
\begin{center}
	ENRICHED $\lambda$-CALCULUS\\[1in]
\end{center}

\begin{flushleft}
The pure $\lambda$-calculus is powerful enough to compute all
computable functions.

However, it is desirable to add constants, such as the integers,
and new rules to manipulate them:
\end{flushleft}

\begin{center}
$\vdots$\\
(+ 0 0) \red 0\\
(+ 0 1) \red 1\\
(+ 1 1) \red 2\\
$\vdots$\\
(+ 12 7) \red 19\\
$\vdots$
\end{center}
\end{slide}


\begin{slide}{}
\begin{center}
	REDUCTION SYSTEMS
\end{center}

\begin{flushleft}

Reduction systems perform computation through equivalence preserving
transformations.

Occurrences of LHS's of rules (REDEX) are replaced by their RHS's
(CONTRACTUM). 
\end{flushleft}

\begin{center}
$(\lambda x . (+ \ x \ 1)) \ 2 $ \red\\
$(+ \ 2 \ 1)$ \red\\
$3$
\end{center}

\begin{flushleft}
Each reduction is semantically an atomic action.
\end{flushleft}
\end{slide}


\begin{slide}{}
\begin{center}
	NORMAL FORMS\\[1in]
\end{center}

\begin{flushleft}
When there are no more redexes, an expression is in NORMAL FORM.

Not all expressions have a normal form:
\end{flushleft}

\begin{center}
	$(\lambda x . x \ x) \ (\lambda x . x \ x)$\\[1in]
\end{center}

\begin{flushleft}
\underline{Church-Rosser I}\\
If an expression has a normal form, it is unique.

\underline{Church-Rosser II}\\
Always contracting the leftmost redex is guaranteed to
to produce the normal form, if it exists.
\end{flushleft}
\end{slide}



\begin{slide}{}
\begin{center}
The Complete $\lambda\beta$-Calculus\\
What is it?\\[2em]
PARTIAL EVALUATION
\end{center}

\begin{flushleft}
The $\lambda$-calculus IS NOT just functional programming.

The $\lambda$-calculus transforms {\em expressions}.\\[1in]



\underline{Example}:
The body of an abstraction is reduced even if the abstraction is
not applied to an argument.
\end{flushleft}
\begin{center}
	$(\lambda x . (+ \ x \ (* \ 2 \ 3)))$ \red\\

	$(\lambda x . (+ \ x \ 6))$
\end{center}
\end{slide}

\begin{slide}{}
\begin{center}
The Complete $\lambda\beta$-Calculus\\
What is it?\\[2em]
PARTIAL EVALUATION (Cont.)
\end{center}
\vspace*{1.5in}
\begin{flushleft}
Proper handling of unbound variables is extended to include 
all language primitives.
\end{flushleft}
\begin{center}
	($\lambda x$ . (equal? [1 2 3] [1 $x$ 3])) \red\\

	($\lambda x$ . (equal? $x$ 2))
\end{center}
\end{slide}


\begin{slide}{}
\begin{center}
The Complete $\lambda\beta$-Calculus\\
What is it?\\[2em]
INCREMENTAL REDUCTION\\[1in]
\end{center}

\begin{flushleft}
The user can specify how many reductions are to be performed.
\end{flushleft}
\begin {itemize}
	\item allows `stepping' through reduction of expressions
	\item debugging tool
	\item guards against non-terminating reduction sequences
	\item partial execution
\end{itemize}
\end{slide}



\begin{slide}{}
\begin{center}
IMPLEMENTATION ISSUES THAT SHAPED RED2
\end{center}

Reduction Order\\
\hspace*{1.5em} Applicative order, Normal order, or \\ 
\hspace*{1.5em} something else?

Representation of Expressions\\
\hspace*{1.5em} \head eliminate variable name clashes\\
\hspace*{1.5em} \head support de-compilation\\
\hspace*{1.5em} \head support reduction order using a \\
\hspace*{2em} minimum of memory bandwidth

Substitution Mechanism\\
\hspace*{1.5em} Delay via environment, Graph mutation,\\
\hspace*{1.5em} or eliminate substitution?
\end{slide}



\begin{slide}{}
\begin{center}
Representation of Expressions\\[.5em]
ELIMINATE NAME CLASHES\\[2em]
\end{center}

\begin{flushleft}
\underline{PROBLEM}\\
Use of symbolic identifiers is useful
for people, but complicates the mechanical manipulation
of expressions.\\[.5in]
\underline{SOLUTION}\\
Use DeBruijn indices, which are based on the structure of
expressions.\\[1in]
\end{flushleft}
\begin{center}
\( (\lambda x . \lambda y . (x \ (y \ (\lambda z . y z)))) \)

\( (\lambda \lambda . (1 \ (0 \ (\lambda.1 \ 0)))) \)
\end{center}

\end{slide}



\begin{slide}{}
\begin{center}
Representation of Expressions\\[.5em]
SUPPORT DE-COMPILATION\\[2em]
\end{center}
Strings or Graphs?
\vspace*{4in}
%put figure 4.1 here
\begin{flushleft}
If expressions are represented as graphs, an inorder traversal yields 
a normal order reduction sequence.
\end{flushleft}
\end{slide}


\begin{slide}{}
\begin{center}
Representation of Expressions\\[.5em]
SUPPORTING REDUCTION/TRAVERSAL\\[2em]
\end{center}

\begin{flushleft}
Graphs can be represented linearly in memory by instructions which
mirror the objects in the $\lambda$-calculus.
\end{flushleft}
\vspace*{2in}
%put lgr1.tex here
\begin{flushleft}
A more compact representation uses a {\em head bit} to indicate if a single
instruction is the operand of an application.
\end{flushleft}
\vspace*{2in}
%put lgr2.tex here
\end{slide}

%13
\begin{slide}{}
\begin{center}
	COMPILATION EXAMPLES\\[.5in]
\end{center}

$\cal C$ [ $(\lambda x . (+ \ x \ 1))\ 5$] \ =
\begin{tabbing}
0:123\=LAMBDA123\=\kill
0: \> \spine INT \> 5\\
1: \> \spine LAMBDA \> $x$\\
2: \> \spine INT \> 1\\
3: \> \spine VAR \> 0\\
4: \> \head PRIM\_2 \> +
\end{tabbing}



$\cal C$ [ $(\lambda a . \lambda b . (a \ b)) \ (\lambda x . x)$] \ = 
\begin{tabbing}
0:123\=LAMBDA123\=\kill
0: \> \spine APP \> 5\\
1: \> \spine LAMBDA \> $a$\\
2: \> \spine LAMBDA \> $b$\\
3: \> \spine VAR \> 0\\
4: \> \head VAR \> 1\\[1em]
5: \> \spine LAMBDA \> $x$\\
6: \> \head VAR \> 0
\end{tabbing}
\end{slide}


\begin{slide}{}
\begin{center}
RED2 MEMORY MODEL
\end{center}

\vspace*{2in}
%put figure 4.7 here

ADDRESS REGISTERS:
\begin{tabbing}
12\={\em env}:12\=\kill
\>{\em pc} \> Program Counter\\
\>{\em fsp} \> Free Space Pointer\\
\>{\em env} \> Environment Pointer\\
\>{\em c} \> Control Stack
\end{tabbing}

DATA REGISTERS:
\begin{tabbing}
12\={\em env}:12\=\kill
\>{\em d} \> Direction Flag\\
\>{\em q} \> Quantum\\
\>$\phi$ \> Binding Offset\\[1em]
%12\={\em argcnt}:12\=\kill
%\>{\em argcnt} \> argument count\\
%\>{\em prim} \> primitive operator\\
%\>{\em fire} \> operator argument count
\end{tabbing}
\end{slide}


\begin{slide}{}
\begin{center}
EXECUTION MODEL
\end{center}
\vspace*{\fill}
\end{slide}

\begin{slide}{}
\begin{center}
Instruction Execution Example\\[.5em]
CONTRACTING A $\beta$-REDEX\\[2em] 
%put lambda.tex picture here
\end{center}
\vspace*{\fill}
\end{slide}


\begin{slide}{}
\begin{center}
Instruction Execution Example\\[.5em]
BINARY ARITHMETIC PRIMITIVE\\[2em]
%put figure 5.1 here
\end{center}
\vspace*{\fill}
\end{slide}


\begin{slide}{}
\begin{center}
ASSESSMENT OF RED2\\[2in]
\end{center}

\begin{itemize}
\item Efficient for the $\lambda$-calculus
\item Compares well with functional language implementations
\item Performs poorly with regard to sharing
\item Suitable for implementation using existing von-Neumann technology
\end{itemize}

\end{slide}


\begin{slide}{}
\begin{center}
SHARING
\end{center}

Memory management policy of RED2 is not conducive to sharing in
general.\\

\vspace*{2in}
%put in Figure 5.2 here

Arguments that reduce to atomic data values can be shared via the
environment. 

Sharing of structures could be done by adding a separate memory area
which is garbage collected.
\end{slide}


\begin{slide}{}
\begin{center}
ADDING LOGIC PROGRAMMING\\[.5in]
\end{center}

\begin{itemize}
\item Relations represented \`{a} la Clark
\item Depth-first SLD-Resolution with left-most selection rule
\item Single solution
\item {\tt UNIFY} primitive
\item Use {\tt AND} and {\tt OR} to control search
\end{itemize}

RED2 Not Well Suited For Logic Programming

\end{slide}
